\paragraph{Static cost~\cite[Lemmas 6.5--6.6, Prop.~6.7 and Cor.~6.8]{tc}.} Let \(\widehat W_e=p_e(1-p_e)\),
and let \(\mu_{\widehat W}\) be the law on permutations supported on the mask, with weights
\(\prod_{e\in\sigma}\widehat W_e\). For an interior stationary point with unit axial masses and positive
slacks, \cite[Lemma~6.5]{tc} gives the identity
\[
 Q(B)e^{-H_B(p)}=\frac{\operatorname{per}\widehat W}{\prod_e(1-p_e)}
                  \mathbb E_{\mu_{\widehat W}}\mathcal A,\qquad
 \mathcal A(\sigma)=
 \frac{\prod_f h_f}{\prod_e(1-p_e)^2}\,\mathbf1_{\{\sigma\ {\rm legal}\}}
 \prod_{e\in\sigma}\frac{(1-p_e)^2}{h_{D(e)}h_{A(e)}} .
\]
Here \(\Gamma_4,\Gamma_2\) denote the canonical queen and rook normalizers with the weights \(p\) of the
same board. If \(Q(B)>0\), \(p_e\le C/m\), and \(\|p-\mathbf1\mathbf1^{\mathsf T}/m\|_{\rm
op}\le\rho_*<1\), one obtains \(\log Q(B)-L(B)=\log\mathbb E_{\mu_{\widehat W}}\mathcal A
-\log(\Gamma_4/\Gamma_2)+o(1)\), uniformly for fixed \(C,\rho_*\). The axial matrix has spectrum
\(1\pm\sigma_i(p)\), so \eqref{eq:clase-importada} implies the last hypothesis.

With the margins of~\eqref{eq:clase-importada}, \cite[Lemma~6.6]{tc} gives
\((S(t)^+)_{uv}=\mathbf1_{\{u=v\}}+O_{C,\eta,\gamma}(m^{-1})\), uniformly in the axial indices and in
\(0\le t\le1\). Here \(S(t)=(1-t)S_{R,2}+tS_{R,4}\), where \(S_{R,d}\) is the axial incidence matrix
weighted by \(p(1-p)/(1-dp)\). With \(\max_f\ell_f\le\epsilon_{\rm ext}<1\), \cite[Prop.~6.7]{tc} uses this
inverse to give \(\log(\Gamma_4/\Gamma_2)=O_{C,\eta,\gamma}(\epsilon_{\rm ext}+m^{-1})\). At a first
mesoscopic exit, \cite[Cor.~6.8]{tc} yields \((\log Q-L)^+\le C_0+\beta m\), where the left-hand side is
read as zero when \(Q=0\), uniformly in the outer class. Its \emph{proof} applies at every arrival in that
class. Indeed, \(\sum_f\log h_f-2\sum_e\log(1-p_e)\le O_C(1)\), since \(\sum_f\ell_f=2m=2\sum_e p_e\) and
\(\sum_ep_e^2\le C\). Each of the \(m\) chosen squares contributes at most \(2\epsilon_{\rm
ext}/(1-\epsilon_{\rm ext})\) to the logarithm of the last product. \cite[Lemmas~6.5--6.6 and Prop.~6.7]{tc}
absorb the remainder into \(C_0\), independently of \(n/m\). Appendix~\ref{app:signed} applies the corollary
in the bands, and Appendix~\ref{app:arrival} applies the static argument at regular arrival with a long
line.
